\section{The seed-request theory and the exact numerical calculus}
\label{sec:seed-theory}

The successful change uses the \emph{seed-request} form of the K schema.
The request at the root then repeats the request already used in the
closure clause.  This makes the expanded source longer but gives its
binary sequence grammar more reusable pieces.  The seed reformulation
and its logical strength already occur in the pinned K230 proof:
\code{KPSeed.v} defines \code{seed_root}, and \code{KPrimeCore.v} defines
\code{seeded_core}.  The contribution here is its use as the actual
emitted formula, together with a grammar and decoder layout that turn
that repetition into a state saving.  We do not claim a new set theory
or a new discovery of the seed equivalence~\cite{kpseed,kprimecore}.

\subsection{An ordinary first-order presentation}

Use classical first-order logic with genuine equality, a nonempty domain,
and one binary relation $R$.  A matrix $\varphi(y,p,\bar q)$ has an output
role $y$, an input role $p$, and any finite tuple of parameters $\bar q$.
All displayed bound variables are chosen capture-free.  Put
\begin{align}
 \Sub(p,u)&:=\forall z\bigl(R(z,p)\to R(z,u)\bigr),
 \label{eq:ks-sub}\\
 W_\varphi(p,C)&:=
   \bigl(\exists y\,\varphi(y,p,\bar q)\bigr)
   \to\exists y\bigl(\varphi(y,p,\bar q)\land R(y,C)\bigr),
 \label{eq:ks-request}\\
 H_\varphi(C)&:=\forall u\Bigl(R(u,C)\to
   \Bigl[\bigl(\exists p\,\varphi(u,p,\bar q)\bigr)
   \land\forall p\bigl(\Sub(p,u)\to W_\varphi(p,C)\bigr)\Bigr]\Bigr).
 \label{eq:ks-closure}
\end{align}
The first conjunct in the square brackets is \emph{support}: every member
of $C$ is an output for some input.  The second conjunct requests witnesses
for subsets of members of $C$.  The support witness is not an additional
parameter of the later requests; their input quantifier binds a new role.

Let $\mathsf K_{\mathrm s}$ contain equality logic and the universal
closures, over $\bar q$, of all instances
\begin{equation}
  \forall a\,\exists C
       \bigl(W_\varphi(a,C)\land H_\varphi(C)\bigr).
  \label{eq:ks-schema}
\end{equation}
Let $\mathsf K_{\mathrm o}$ denote the inherited output-root theory,
whose corresponding instance is
\begin{equation}
 \forall a\,\forall x\,\exists C
 \bigl[(\varphi(x,a,\bar q)\to R(x,C))\land H_\varphi(C)\bigr].
 \label{eq:ko-schema}
\end{equation}
The presentation of $H$ in~\eqref{eq:ks-closure} is classically equivalent
to the exact upstream closure.  In that closure a support existential
has the form $\exists p(\varphi(u,p,\bar q)\land J)$, and the first
quantifier of $J$ rebinds $p$.  Thus $p$ is not free in $J$, and the
existential can be moved into the first conjunct.  The upstream witness
request uses
\begin{equation}
 \forall y\Bigl(\varphi(y,p,\bar q)\to
       \exists y\bigl(\varphi(y,p,\bar q)\land R(y,C)\bigr)\Bigr).
 \label{eq:ks-universal-request}
\end{equation}
The existential on the right has no free $y$, so this is classically
equivalent to~\eqref{eq:ks-request}.  These are finite first-order proof
transformations, with the variable conditions stated explicitly.

\begin{theorem}[Calibration of the seed-request theory]
\label{thm:ks-calibration}
With $\IR:=\forall x\,\neg R(x,x)$, ZFC proves
\[
 \Con(\mathsf K_{\mathrm s})
 \quad\longleftrightarrow\quad
 \Con(\mathsf K_{\mathrm s}+\IR)
 \quad\longleftrightarrow\quad
 \Con(\mathrm{ZFC}).
\]
More precisely, the two inconsistency directions are effective
transformations of finite derivations, formalizable in ZFC.
\end{theorem}

\begin{proof}
First, $\mathsf K_{\mathrm o}$ proves every seed-request instance.  Fix
$\varphi,a,\bar q$ and split classically on
$\exists x\,\varphi(x,a,\bar q)$.  If an output $x$ exists, apply the
output-root instance to that $x$.  Its root puts $x$ into $C$, hence gives
$W_\varphi(a,C)$, and its closure gives $H_\varphi(C)$.  If there is no
output, apply the output-root instance to any domain element; the
resulting $C$ still has $H_\varphi(C)$, and its root request is vacuous.
This is the argument formalized by \code{output_root_to_seed} in
\code{KPSeed.v}, followed by \code{kprime_core_to_seed} in
\code{KPrimeCore.v}~\cite{kpseed,kprimecore}.

The pinned upper interpretation realizes $\mathsf K_{\mathrm o}+\IR$
in the constructible universe.  In finite-proof form, the theorem
\code{KPrimeOrdIR_standardZFC_inconsistent} in
\code{KUpperConsistency.v} transforms each refutation of that theory
into a ZFC refutation~\cite{k230}.  Replacing seed axioms by their
output-root derivations therefore gives
\[
 \Con(\mathrm{ZFC})\quad\Longrightarrow\quad
 \Con(\mathsf K_{\mathrm s}+\IR)\quad\Longrightarrow\quad
 \Con(\mathsf K_{\mathrm s}).
\]
This use of a definable interpretation does not require a set model of
all of ZFC inside ZFC.

For the converse, consider the factored theory $\mathsf K_0$ from the
same development.  Its schema asks for a $C$ satisfying support and
\begin{equation}
 \forall p\Bigl(
   \bigl[p=a\ \lor\ \exists u(R(u,C)\land\Sub(p,u))\bigr]
       \to W_\varphi(p,C)\Bigr).
 \label{eq:ks-k0-request}
\end{equation}
Every witness $C$ to~\eqref{eq:ks-schema} satisfies these requirements.
If $p=a$, equality transports the root request to $p$.  If
$u\in C$ and $\Sub(p,u)$, the closure clause at $u$ supplies the request
at $p$.  The first conjunct of that same closure supplies support.
This is the finite derivation \code{seeded_core_to_kp0} in
\code{KPrimeCore.v}, with only the displayed conjunction and request
conversions~\cite{kprimecore}.

The upstream quotient interpretation already factors through
$\mathsf K_0$: \code{K0_ZFc_inconsistent} in
\code{KQuotZFTransfer.v} transforms a refutation of its ZF base into a
$\mathsf K_0$ refutation.  Composing it with the Foundation and Choice
relative-consistency reduction \code{ac_reg_relcon}, as in
\code{KQuotZFC.v}, gives
\[
 \mathrm{ZFC}\vdash\bot
 \quad\Longrightarrow\quad
 \mathsf K_0\vdash\bot
 \quad\Longrightarrow\quad
 \mathsf K_{\mathrm s}\vdash\bot.
\]
Here the first and last occurrences of $\vdash$ describe finite
object-theory derivations; the implication between them is the
metatheoretic proof transformation~\cite{kquot}.
Adding $\IR$ preserves this direction.

Each additional conversion above is a fixed finite logical template
applied to an input matrix.  It composes with the inherited finite-proof
transformations.  Coding these finite constructions and verifying them
by induction in ZFC proves the stated internal equivalences.  This last
claim is an ordinary mathematical formalizability argument, not an
identification of Rocq's metatheory with ZFC.
\end{proof}

\subsection{The expression actually emitted}
\label{sec:literal-ks}

The native constructors are
\begin{align*}
 P(x,y)&=\frac{(x+y)(x+y+1)}2+y,\\
 \cI(x,y)&=4P(x,y), & \cR(x,y)&=4P(x,y)+1,\\
 \cQ(x,y)&=4P(x,y)+2, & \cE(x,y)&=4P(x,y)+3.
\end{align*}
The first argument of $\cQ$ is a variable label.  The other constructors
represent implication, membership and primitive equality.  Write
$\cF=1=\cR(0,0)$ and $\cN(t)=\cI(t,1)$.  For reference, the unfused
negative encodings of conjunction and existence are
\[
 A(x,y)=\cN\bigl(\cI(x,\cN(y))\bigr),\qquad
 X_v(x)=\cN\bigl(\cQ(v,\cN(x))\bigr).
\]
These equations specify natural-number operations; an equation between
these codes is stronger than a logical equivalence of decoded formulas.

The source retains the eight logical rows $\rho_1,\ldots,\rho_8$, the
truth seed $0$, modus ponens, theorem generalization, and the gate
$c\ne a,b$ from K224.  In particular it retains the restricted equality
helper, rather than replacing it by unrestricted equality vacuity.
The classical propositional derivations used below are available in
this numerical core~\cite{klogic,k224}.  One useful detail is that
relative falsity is sufficient: the retained $\rho_6$ has the form
$X_v(\cQ(v,t))\to t$.  Since $X_v(\cQ(v,t))$ has the form $U\to\cF$,
the derived implicational calculus gives
$\cF\to X_v(\cQ(v,t))$ and hence $\cF\to t$.  Thus double-negation
elimination and the negative conjunction rules do not require a new
logical row or a freshness assumption about label $0$.

For each enabled triple $(a,b,c)$, define the following exact arithmetic
expressions, with $a$ the input label and $b$ the output label:
\begin{align}
 r&=\cR(b,c),\nonumber\\
 s&=\cQ\bigl(c,\cI(\cR(c,a),\cR(c,b))\bigr),\nonumber\\
 w&=\cQ\bigl(b,\cI(c,\cN(\cQ(b,\cI(c,\cN(r)))))\bigr),
 \label{eq:literal-w}\\
 u&=\cQ\bigl(a,\cI(s,w)\bigr),\nonumber\\
 h&=\cQ\bigl(b,\cI(r,\cN(\cQ(a,\cI(c,\cN(u)))))\bigr),
 \label{eq:literal-h}\\
 v&=\cI\bigl(\cR(a,b),\cQ(c,\cR(a,b))\bigr),\nonumber\\
 e&=\cI\bigl(\cE(a,b),\cQ(\cR(a,b),\cE(a,b))\bigr).
 \label{eq:literal-v-e}
\end{align}
The actual released bundle is
\begin{equation}
 B_{223}(a,b,c)=
 \cN\Bigl(\cI\Bigl(v,\cI\bigl(e,
                   \cQ(c,\cI(h,\cN(w)))\bigr)\Bigr)\Bigr).
 \label{eq:literal-bundle}
\end{equation}
For a disabled triple, the emitter returns $0$.

In particular, the matrix is the literal natural number $c$, and the
private container binder is the same literal number $c$.  It has not
been replaced by the code of a more convenient decoded formula.
The decoded content of~\eqref{eq:literal-bundle} is classically
\begin{equation}
 V_R(a,b,c)\land
 \bigl(H_R(a,b)\land\exists c\,(H\land W)\bigr),
 \label{eq:literal-bundle-content}
\end{equation}
where $V_R$ is membership vacuity, $H_R$ is the inherited equality helper,
and $H,W$ are the named forms of~\eqref{eq:ks-closure}
and~\eqref{eq:ks-universal-request}.  Their grouping and order matter
for the source grammar even though they do not matter for classical
provability.  The numerical core in isolation would be
$\cN(\cQ(c,\cI(h,\cN(w))))$.  Placing it last in the bundle contributes
one more outer negation; the resulting double negation cancels to the
$\cQ(c,\cI(h,\cN(w)))$ visible in~\eqref{eq:literal-bundle}.
This cancellation crosses the boundary between the reordered bundle and
its seed core.

\subsection{Decoding every enabled triple, including low codes}
\label{sec:ks-decoding}

The inherited total decoder $D$ treats two numerical cases specially:
\[
 D(0)=\bigl(\forall 0\,(0=0)\bigr)
              \to\bigl(\forall 0\,(0=0)\bigr),
 \qquad D(1)=R(0,0).
\]
For any $t$, it decodes $\cI(t,1)$ as genuine named negation of $D(t)$.
For other implications, the exact theorem is
\[
 D(\cI(x,y))=D(x)\to D(y)
 \quad\text{if }y\ne1\text{ and }(x,y)\ne(0,0).
\]
The atomic and universal-quantifier decoding identities are unconditional.
These are \code{D_0}, \code{D_wn}, \code{D_wim}, \code{D_wel},
\code{D_weq} and \code{D_wal} in \code{KDecoder230.v}~\cite{kdecode}.

\begin{lemma}[Exact constructor interface]
\label{lem:ks-decode-interface}
For all naturals $a,b,c$ with $c\ne a,b$, applying $D$ to
\eqref{eq:literal-bundle} gives exactly the stated named formula with
matrix $D(c)$, before the ordinary classical conversions in
\eqref{eq:literal-bundle-content}.  No non-negation implication is
misread as a negation or as code $0$, and $B_{223}(a,b,c)>1$.
\end{lemma}

\begin{proof}
Each intentional $\cN$ has literal consequent $1$, so use \code{D_wn}.
At each other implication the consequent is a quantifier, a negation,
or another positive compound code, except for the displayed membership
consequents.  The potentially exceptional membership consequents are
$\cR(c,b)$ and $\cR(b,c)$; the gate excludes $b=c$, so neither is
$\cR(0,0)=1$.  The other cases follow immediately from their residues
modulo $4$.  In each non-negation implication at least one argument is
positive, ruling out $(0,0)$.  Induction over the finite displayed tree
therefore applies the decoding identities at every constructor.
The outer negation has a positive code divisible by $4$, hence exceeds
$1$.  The included decoder-side-condition checker verifies this same
finite, uniform constructor argument on the exact expression tree; it
does not enumerate a bounded range of triples.
\end{proof}

The essential freshness theorem is
\begin{equation}
             c\notin\operatorname{FV}(D(c))
             \quad\text{for every natural }c.
 \label{eq:private-freshness}
\end{equation}
For $c>0$, every variable name, free or bound, occurring in $D(c)$ is
strictly below $c$.  At $c=0$, the stronger assertion about all occurrences
would be false: the displayed tautology contains bound $0$'s.  It is
closed, which is precisely the property needed.  At $c=1$, the only name
in $D(c)=R(0,0)$ is $0<1$.  These facts are proved by
\code{D_names_below}, \code{D_private_fresh}, \code{D_zero_closed} and
\code{D_one} in \code{KDecodeFresh.v}~\cite{kdecode}.

If $a\ne b$, abstract the free named occurrences of $b$ as the output,
those of $a$ as the input, and the finitely many other free names as
parameters.  The environment equalities in \code{KDecodeAbstract.v}
then identify the membership, subset, input rebinding, output rebinding
and matrix clauses with the ordinary seed-request instance.  Their
hypotheses are exactly $a\ne b$, $c\ne a,b$ and
\eqref{eq:private-freshness}.  Thus this argument applies to every
$D(c)$, including $D(0)$ and $D(1)$, rather than only to padded codes
chosen for a later completeness argument~\cite{kdecode}.

\begin{lemma}[Coincident active labels]
\label{lem:ks-diagonal}
When $a=b$ and $c\ne a$, the decoded seed core is valid in the upper
set-theoretic interpretation.
\end{lemma}

\begin{proof}
Fix the other free names and write $\psi(x)$ for $D(c)$ evaluated with
its single active name $a=b$ assigned $x$.  The private label $c$ stays
fresh.  Rename locally bound variables for readability.  The request
becomes
\[
 W(C)=\bigl(\exists x\,\psi(x)\bigr)
               \to\exists x\bigl(\psi(x)\land R(x,C)\bigr).
\]
In the subset test both endpoints now use the same active name, so it
is $\forall z(R(z,p)\to R(z,p))$.  The support existential and subsequent
request quantifiers rebind that same active name and shadow the member
variable of the outer guard.  The remaining universal request is therefore
$\forall p(\Sub(p,p)\to W(C))$; $W(C)$ has no free $p$, so this is
equivalent to $W(C)$ in a nonempty domain.  Consequently the exact closure
simplifies to
\[
 H(C)=\forall u\Bigl(R(u,C)\to
                 \bigl[(\exists x\,\psi(x))\land W(C)\bigr]\Bigr).
\]
If $\neg\exists x\,\psi(x)$, choose an empty $C$: its request is vacuous
and its closure has no members to check.  If a witness $x$ exists, choose
$C=\{x\}$.  Then the same true existence statement and the same $W(C)$
verify every guarded closure clause.  Empty sets and singletons are
available in the constructible interpretation used for the upper
consistency direction.  The case split is on existence of a witness,
not on the matrix evaluated at the old value of the active variable.
No two-distinct-role abstraction has been used in this case.
\end{proof}

Genuine named negation is used before the preceding first-order
simplifications.  For example, at output label $0$ one must not infer
freshness by expanding existence as an implication to the raw atom
$R(0,0)$.  The correct decoding first uses \code{D_wn} and
\code{D_wex}; the inherited relative-falsity/$\IR$ translation handles
the numerical proof system afterwards.  Its exceptional-code and
modus-ponens cases are unchanged~\cite{kdecode,klogic,k224}.

\subsection{Why the inherited equality helper is sufficient}
\label{sec:ks-equality}

The two helper formulas in the literal bundle are
\begin{align*}
 V_R(a,b,c)&:\quad R(a,b)\to\forall c\,R(a,b),\quad c\ne a,b,\\
 H_R(a,b)&:\quad (a=b)\to
                    \forall_{\kappa(a,b)}(a=b),
 \qquad \kappa(a,b)=\cR(a,b).
\end{align*}
Because $P(a,b)\ge\max(a,b)$, one has
$\kappa(a,b)>\max(a,b)$, including $(a,b)=(0,0)$.
Thus every $H_R$ is sound with genuine equality.  Every pair receives
its helper in the numerical theory: choose any original bundle parameter
$c\ne a,b$, and project the helper using the retained conjunction rules.
Its binder is $\kappa(a,b)$ independently of that choice of $c$.

For the lower consistency direction, use the inherited interpretation of
equality by membership indiscernibility.  For each pair of names choose
$z=\cE(a,b)>a,b$ and let
\begin{equation}
 E_{\!R}(a,b):=
 \bigl[\forall z\,(R(z,a)\leftrightarrow R(z,b))\bigr]\land
 \bigl[\forall z\,(R(a,z)\leftrightarrow R(b,z))\bigr].
 \label{eq:ks-defined-equality}
\end{equation}
The biconditionals and conjunction are expanded in the relative-falsity
syntax.  Let $\theta$ replace primitive equality by this fixed formula,
replace genuine bottom by $\cF=R(0,0)$, and preserve membership,
implication and each universal binder.

The logical calibration restored in K224 from ProveIt Source~116 gives
finite derivations of the translations of all full first-order equality
axioms.  It derives membership-only fresh vacuity by induction, proves
safe-auxiliary transport for~\eqref{eq:ks-defined-equality}, and supplies
interpreted reflexivity, symmetry, transitivity, membership congruence,
equality existence and vacuity for all coincidences of names.  The
helper first supplies the off-diagonal bridge
$(a=b)\to E_{\!R}(a,b)$; on the diagonal the defining membership
identities already prove $E_{\!R}(a,a)$.  These are precisely the logical
parts of Source~116 and the K224 helper calibration.  They use the eight
retained logical rows, membership vacuity, $H_R$, modus ponens and
theorem generalization, independently of the choice of set-theoretic
schema~\cite{source116,k224}.

This gives a proof interpretation into the smaller primitive-equality
calculus.  It does not assert that all original \emph{primitive}
equality-vacuity instances follow from $H_R$.  Nor does it require
$\theta$ to commute literally with arbitrary substitution or alpha
conversion: changing an equality's endpoints can change its numerical
auxiliary $z$.  The explicit translated axiom derivations and auxiliary
transport proofs are used at those steps.

\subsection{Choosing the private label after interpreting the matrix}
\label{sec:ks-code-coupling}

The literal coupling of matrix code and binder is the main additional
obligation in transferring an ordinary schema to this numerical theory.
The equality interpretation usually changes the matrix's code, so
choosing a private label from the old code would not establish the
required source instance.

Take a finite named representative $f$ of an arbitrary ordinary
seed-request matrix.  Use input name $1$, output name $3$, distinct
positive parameter names, and capture-free bound names, as in the
pinned ordinary-to-named construction.  \emph{First} translate equality
and then pad the translated matrix:
\begin{equation}
       \psi=\theta(f),\qquad
       p^{+}=\neg\neg\psi,\qquad
       c=\enc(p^{+}).
 \label{eq:ks-code-choice}
\end{equation}
Here $\enc$ is the actual homomorphic native encoder, and the negations
on the right are encoded using $\cN$.  The elementary code-label
induction gives $c>3$, and every name in $p^{+}$ is strictly below $c$.
In particular $c$ is fresh after all membership-indiscernibility
auxiliaries have been inserted.  It is also fresh in $f$: the translation
retains membership endpoints, equality endpoints and old binder names.

Now choose the \emph{ordinary comparison instance} whose matrix is
$\neg\neg f$ and whose private container name is this $c$.  Ordinary
first-order K allows any fresh private name.  Double-negation equivalence,
capture-free naming and universal closure recover the desired ordinary
seed instance from this selected one.  The outer seed template contains
only membership, implication, negation, quantifiers and matrix
insertions.  It therefore commutes literally with $\theta$, once the
matrix and private binder have been selected:
\begin{equation}
 \theta\bigl(\mathcal K_{\mathrm s}(1,3,c;\neg\neg f)\bigr)
   =\mathcal K_{\mathrm s}(1,3,c;\neg\neg\theta(f)).
 \label{eq:ks-theta-commutes}
\end{equation}
The matrix on the right has native code \emph{exactly} $c$.
The native encoding of that right-hand formula is consequently the
unfused seed template at the actual source parameters $(1,3,c)$.
The gate is enabled.  Projecting the seed component of
$B_{223}(1,3,c)$ and using the retained finite double-negation and
conjunction derivations gives that unfused template in its required
ordinary grouping.  No equality between fused and unfused numerical
codes is being asserted.

This order of construction removes a possible circularity.  The matrix
is translated before its code is chosen; the ordinary comparison
instance's private name is then chosen from that translated code.  There
is no independent matrix integer and no subsequent replacement of the
actual private label.  The argument extends K224's same coupling method
from the output-root schema to the seed-request template~\cite{k224}.

\begin{theorem}[Calibration of the literal numerical theory]
\label{thm:ks-native-calibration}
Let $\mathcal N_{223}$ be the native numerical calculus with the retained
logical core and with the gated bundle~\eqref{eq:literal-bundle}.
ZFC proves
\[
 \bigl(\mathcal N_{223}\text{ has a finite derivation of code }1\bigr)
 \quad\longleftrightarrow\quad
 \bigl(\mathrm{ZFC}\text{ has a finite refutation}\bigr).
\]
Thus ZFC proves
$\Con(\mathcal N_{223})\leftrightarrow\Con(\mathrm{ZFC})$.
\end{theorem}

\begin{proof}
For soundness, decode a finite numerical derivation with the pinned
relative-falsity/$\IR$ proof translation.  Logical rows, truth code $0$,
modus ponens and generalization use the unchanged inherited argument.
Membership vacuity and $H_R$ are sound with genuine equality.  Every
new core is interpreted in ZFC: for distinct active names use the upper
interpretation constructed in the proof of
Theorem~\ref{thm:ks-calibration}, together with the exact abstraction above; for
coincident active names use Lemma~\ref{lem:ks-diagonal}.  The decoded
conclusion is $R(0,0)$, contrary to the interpreted $\IR$.
This yields a finite ZFC refutation from a finite numerical one.

For completeness, a finite ZFC refutation gives a finite
$\mathsf K_{\mathrm s}$ refutation by
Theorem~\ref{thm:ks-calibration}.  Replace its finitely many schema
instances by the ordinary comparison instances selected using
\eqref{eq:ks-code-choice}.  Translate the resulting finite logical
proof by $\theta$.  Section~\ref{sec:ks-equality} supplies every
translated logical axiom, and~\eqref{eq:ks-theta-commutes} together with
projection from the actual enabled bundles supplies every translated
schema instance.  The translation preserves modus ponens and theorem
generalization, and takes the conclusion to $\cF=1$.
If the comparison proof is formulated with $\IR$ as well, choose its
closed named instance $\forall0\,\neg R(0,0)$; its translation is the
generalized identity $\forall0(\cF\to\cF)$.

All selections in this argument are effective operations on finite
syntax: finite maxima choose names, finite recursion translates and
encodes formulas, and fixed logical proof templates perform the
conversions.  The inherited reductions and these additional inductions
can therefore be expressed and verified in ZFC.  The claim concerns
finite proofs and holds internally under ordinary arithmetization; it
is not an inference from tests on a finite collection of matrix codes.
\end{proof}

\subsection{Independent binding to the released source}
\label{sec:ks-emitter-binding}

The preceding theorem is a statement about the \emph{specified}
numerical theory.  Its connection to the released register source is
also checked independently.  The audit
\code{audit_ks223_emitter.py} executes the recorded postfix word
symbolically and compares the resulting tree with a separately written
instance of~\eqref{eq:literal-bundle}.  The selected word has exactly
$77$ macro invocations, using these eight pinned K230 macros:
\begin{center}
\begin{tabular}{@{}rl@{}}
\toprule
\textbf{Dag230 node} & \textbf{Stack operation}\\
\midrule
166 & Push the input label $a$ from register $4$\\
121 & Push the output label $b$ from register $5$\\
168 & Push the actual matrix/private value $c$ from register $3$\\
119 & Apply $\cI$ to the top two operands\\
181 & Apply $\cR$ to the top two operands\\
180 & Apply $\cQ$ to the top two operands\\
182 & Apply $\cE$ to the top two operands\\
218 & Apply $\cN$ to the top operand\\
\bottomrule
\end{tabular}
\end{center}
Their contracts are the explicit \code{par166}, \code{par121},
\code{par168}, \code{cons119}, \code{wel181}, \code{wal180},
\code{weq182} and \code{wn218} lemmas of \code{ProgLib230.v} at the
pinned upstream commit~\cite{k230}.  The symbolic execution preserves
the old stack beneath the one emitted result.

The audit compares each optimized macro with its pinned raw source by
finite bisimulation for all natural register vectors.  It then constructs
an independent right-associated sequence consisting of the six unchanged
closed prefix blocks, the $77$ macros, and the two unchanged closed suffix
blocks.  Its $4\,167$ primitive instruction occurrences match the released
closed emitter region exactly, and the source outside that region is
unchanged.  A second all-register bisimulation relates that canonical
sequence to the released parenthesization.  Its complete certificate has
$1\,240$ boundary pairs.  At related decrement boundaries both programs
test the same register; both possible successors are related; and the
finite intervening increment vectors agree.  Induction on these segments
establishes the register contract for all executions, including
nontermination.  This is a finite control-graph proof, not a runtime
sample or a bounded counter test.

The independent receipt also checks
Lemma~\ref{lem:ks-decode-interface} against the exact symbolic tree.
It records the source and grammar hashes, the $77$-macro word, all eight
macro comparisons, and the complete emitter bisimulation.  The separate
machine sections connect this source contract to exhaustive proof search,
blank-tape initialization and the unbounded binary runtime.  Together
those links apply Theorem~\ref{thm:ks-native-calibration} to the literal
223-state transition table.  No newly compiled end-to-end Rocq or Lean
endpoint for that table is claimed by this report.
